\documentclass[12pt]{article}
\usepackage{amsmath,amssymb}

\begin{document}
\section*{{2022 AMC 12B Problems/Problem 11}}
Source: AoPS Wiki

\subsection*{Solution}
Converting both summands to exponential form, \begin{align*} -1 + i\sqrt{3} &= 2e^{\frac{2\pi i}{3}}, \\ -1 - i\sqrt{3} &= 2e^{-\frac{2\pi i}{3}} = 2e^{\frac{4\pi i}{3}}. \end{align*} Notice that the two terms in the problem are two of the third roots of unity (that is, both of them equal $1$ when raised to the power of $3$ ).
When we replace the summands with their exponential form, we get \[f(n) = \left(e^{\frac{2\pi i}{3}}\right)^n + \left(e^{\frac{4\pi i}{3}}\right)^n.\] When we substitute $n = 2022$ , we get \[f(2022) = \left(e^{\frac{2\pi i}{3}}\right)^{2022} + \left(e^{\frac{4\pi i}{3}}\right)^{2022}.\] We can rewrite $2022$ as $3 \cdot 674$ , how does that help? \[f(2022) = \left(e^{\frac{2\pi i}{3}}\right)^{3 \cdot 674} + \left(e^{\frac{4\pi i}{3}}\right)^{3 \cdot 674} = \left(\left(e^{\frac{2\pi i}{3}}\right)^{3}\right)^{674} + \left(\left(e^{\frac{4\pi i}{3}}\right)^{3}\right)^{674} = 1^{674} + 1^{674} = \boxed{\textbf{(E)} \ 2}.\] Since any third root of unity must cube to $1$ . ~ zoomanTV

\end{document}
